\section{Proof of the main theorem}
\label{sec:main}

We now assemble the constructions of \S\S4--7 with the perturbative
nearly--spherical estimate of \S8. Throughout, $\Om\subset\R^2$ is bounded
open with $\abs\Om=\pi$, $u\in H^1_0(\Om)$ is its torsion function, $\dT=\dT(\Om)=
\tfrac\pi8-T(\Om)\ge0$ is the Saint--Venant deficit, and $\dstar(\Om)=
\tfrac8\pi\dT$ is the scale--invariant relative torsion deficit. We write
$C$ for an absolute constant whose value may change from line to line, and
$\eta_G,C_G$ for the constants of the radial--graph theorem
\cref{thm:graphG}, which depend only on the absolute mode-one constant $K$ of
\eqref{eq:mode1-scoped} below and are therefore themselves absolute. The
threshold $\delta_\star>0$ is absolute and may be shrunk finitely many times.

The section has two results. \Cref{thm:reduction} reduces a near--extremal
$\Om$ to a measurable small--amplitude radial graph $G$ that inherits the
torsion deficit and is close to $\Om$ in measure; this is the entire payload
of \S\S4--7. \Cref{thm:main} then feeds $G$ into \cref{thm:graphG} and
transfers the resulting asymmetry bound back to $\Om$ by the triangle
inequality.

\subsection{The reduction}

The construction selects an inset level $\hat v\sim\sqrt\dT$, passes to the
trapping annulus, cuts the radial folds to make an honest graph, and
optimises the free scale $\tau$ at $\tau=\sqrt\dT$. We carry $\tau$ symbolic
through the estimates and set $\tau=\sqrt\dT$ only at the end, so that the
balance producing the exponent $\tfrac12$ is visible.

\begin{theorem}[Sharp reduction to a graph inset]\label{thm:reduction}
There is an absolute $\delta_\star>0$ such that the following holds. Let
$\Om\subset\R^2$ be bounded open with $\abs\Om=\pi$ and $\dT\le\delta_\star$,
and put
$\tau:=\sqrt\dT$. Then there exist a point $z\in\R^2$ and a bounded open
star-shaped radial graph
\[
   G=\{z+r\omega:\ 0\le r<\rho(\omega),\ \omega\in S^1\},
   \qquad \rho\colon S^1\to(0,\infty)\ \text{measurable},
\]
with $\hat r_G:=\sqrt{\abs G/\pi}$, such that
\begin{enumerate}[label=\textup{(\roman*)},leftmargin=2.6em]
\item \textup{(small amplitude)}
   $\ \norm{\dfrac{\rho}{\hat r_G}-1}_{L^\infty(S^1)}\ \le\ C\,\dT^{1/8}
   \ \le\ \eta_G$;
\item \textup{(controlled deficit)} $\ \dstar(G)\ \le\ C\,\dstar(\Om)$;
\item \textup{(closeness)} $\ \abs{\Om\triangle G}\ \le\ C\,\sqrt\dT$.
\end{enumerate}
Every constant is absolute. The construction uses only the two level
deficits $D_H,D_I$ (through \cref{thm:level-identity,lem:profile,lem:muprime,%
thm:annulus,lem:two-flux}); it invokes no selection principle, no transport,
and no Serrin/Magnanini--Poggesi symmetry estimate.
\end{theorem}

\begin{proof}
Throughout we abbreviate $\diso=\diso(\Etil)$, $\Exc=P(\Etil)-2\pi\hat r$,
$a=\abs{\Etil\triangle B_{\hat r}(z)}$ and $\eta=\rho_e-\rho_i$, with the
notation of \cref{def:inset,thm:annulus}.

\emph{Step 0: the inset.} Apply \cref{lem:level} with the given $\tau\in
(0,\tfrac1{16})$ (legitimate since $\tau=\sqrt\dT\le\sqrt{\delta_\star}<
\tfrac1{16}$ once $\delta_\star<\tfrac1{256}$): there is a regular value
$\hat v\in(\tau,2\tau)$ with
\begin{equation}\label{eq:m-good}
   D_H(\hat v)+D_I(\hat v)\le A_0\,\frac\dT\tau,
   \qquad
   \mu_B(\hat v)-\mu(\hat v)\le A_0\,\frac\dT\tau,
\end{equation}
$A_0$ absolute. Let $\Etil_0$ be the connected inset of \cref{def:inset} (the
component of $\{u>\hat v\}$ of largest measure); by \cref{lem:noholes} it is
connected and has no interior holes, and $\mu_E=\abs{\Etil_0}$,
$\hat r=\sqrt{\mu_E/\pi}$. By \cref{lem:close,lem:spill},
\begin{equation}\label{eq:m-close}
   \abs{\Om\triangle\Etil_0}\ \le\ 8\pi\tau+A_0\frac\dT\tau
       +\frac{A_0^2}{16\pi^3}\,\frac{\dT^2}{\tau^2}.
\end{equation}
Since $\mu_E=\abs{\Etil_0}=\pi-\abs{\Om\triangle\Etil_0}$ (as $\Etil_0\subset\Om$),
\eqref{eq:m-close} gives $\mu_E\ge\pi-(8\pi\tau+A_0\dT/\tau+A_0^2\dT^2/
(16\pi^3\tau^2))$; for $\delta_\star$ small (so that the right--hand bound is
$\le\pi/2$, which is exactly the standing hypothesis \eqref{eq:halfmass} of
\cref{sec:inset} at $\tau=\sqrt\dT$, with the higher--order spillover term
absorbed) we obtain $\mu_E\ge\pi/2$, hence $\hat r=\sqrt{\mu_E/\pi}\ge1/\sqrt2$
and, for the annulus produced below, $\rho_i\ge\tfrac12$ after possibly
shrinking $\delta_\star$.

\emph{Step $0'$: hole-filling.} Pass from $\Etil_0$ to its saturation
$\Etil:=E^h$ of \cref{def:fill}, with all bounded complementary components
filled in (\cref{conv:fill}). By \cref{lem:fill}\ref{it:fill-top} the set
$\Etil$ is bounded, open, connected and \emph{globally hole-free} ---
$\R^2\setminus\overline\Etil$ connected --- \emph{unconditionally}, with no
hypothesis on $\partial\Om$; this discharges the exterior-connectedness
convention \cref{lem:ext-conn} required by the trapping annulus and the
inscribed-ball lemma. The filled mass obeys $H:=\abs{\Etil\setminus\Etil_0}
\le D_I(\Etil_0)^2/(64\pi^2\mu_E)\le C(\dT/\tau)^2$
(\cref{lem:fill}\ref{it:fill-key},\ref{it:fill-scales}), so at $\tau=\sqrt\dT$,
$H\le C\dT$, of the same higher order as the spillover. By
\cref{lem:fill}\ref{it:fill-scales} all inset scales are preserved with the
same rates: $\diso(\Etil)\le\diso(\Etil_0)$, $\Exc(\Etil)\le\Exc(\Etil_0)+CH$,
$a(\Etil)\le a(\Etil_0)+CH$, $\dstar(\Etil)\le\dstar(\Etil_0)+CH/\mu_E$, and
the closeness deteriorates only by $H$:
\begin{equation}\label{eq:m-fillclose}
   \abs{\Om\triangle\Etil}\ \le\ \abs{\Om\triangle\Etil_0}+H
   \ \le\ 8\pi\tau+A_0\frac\dT\tau+\frac{A_0^2}{16\pi^3}\frac{\dT^2}{\tau^2}+CH.
\end{equation}
From here on $\Etil$ denotes the hole-filled inset, $\mu_E:=\abs\Etil$ (now
$\mu_E+H$ of the previous step, still $\ge\pi/2$), $\hat r:=\sqrt{\mu_E/\pi}$,
and we run Steps~1--4 on it; all of \cref{sec:annulus,sec:fold,sec:cut} apply
to it directly (\cref{conv:fill}).

\emph{Step 1: the scales (carrying $\tau$ symbolic).} From the inset control
\eqref{eq:diso-inset} of \cref{prop:insetcontrols} (which includes the
spillover term $4\pi S$ and absorbs it),
\begin{equation}\label{eq:m-diso}
   \diso\ \le\ C\,\frac\dT\tau,
   \qquad
   \Exc=2\pi\hat r\,\diso\ \le\ C\,\frac\dT\tau,
\end{equation}
using $\hat r\le1$. The sharp quantitative
isoperimetric inequality \cref{thm:fmp}, in the form
$a=\abs{\Etil\triangle B_{\hat r}(z)}\le C_{\mathrm F}\sqrt{\diso}\,\abs\Etil$,
gives
\begin{equation}\label{eq:m-a}
   a\ \le\ C\,\sqrt{\diso}\ \le\ C\Bigl(\frac\dT\tau\Bigr)^{1/2}.
\end{equation}
The trapping annulus \cref{thm:annulus}, applied to the hole-filled inset
$\Etil=E^h$ (globally hole-free unconditionally by \cref{lem:fill}\ref{it:fill-top},
\cref{conv:fill}), yields
$B_{\rho_i}(z)\subset\Etil\subset B_{\rho_e}(z)$ with
\begin{equation}\label{eq:m-eta}
   \eta=\rho_e-\rho_i\ \le\ C\,(\Exc+a)^{1/2}.
\end{equation}
By \cref{cor:gd} the radial--graph defect of $\Etil$ about $z$ obeys
\begin{equation}\label{eq:m-gd}
   \gd(\Etil)\ \le\ C\,\eta\,(\Exc+a\eta).
\end{equation}

\emph{Step 2: optimise $\tau=\sqrt\dT$ and re--derive every exponent.} Set
$\tau=\sqrt\dT$, so $\dT/\tau=\sqrt\dT$. We re--derive each scale
independently from \eqref{eq:m-diso}--\eqref{eq:m-gd}:
\begin{align}
   \diso\ &\le\ C\,\frac\dT\tau\ =\ C\,\dT^{1/2}, &
   \Exc\ &\le\ C\,\dT^{1/2}, \label{eq:s-diso}\\
   a\ &\le\ C\Bigl(\frac\dT\tau\Bigr)^{1/2}\ =\ C\,\dT^{1/4}, \label{eq:s-a}\\
   \eta\ &\le\ C\,(\Exc+a)^{1/2}\ \le\ C\,a^{1/2}\ \le\ C\,\dT^{1/8},
       \label{eq:s-eta}
\end{align}
where in \eqref{eq:s-eta} we used that $a\sim\dT^{1/4}$ dominates
$\Exc\sim\dT^{1/2}$ as $\dT\to0$, so $\Exc+a\le 2a\le C\dT^{1/4}$ and
$(\Exc+a)^{1/2}\le C\dT^{1/8}$. For the graph defect, \eqref{eq:m-gd} gives
\begin{equation}\label{eq:s-gd}
   \gd(\Etil)\ \le\ C\,\eta\,(\Exc+a\eta)
   \ \le\ C\bigl(\dT^{1/8}\!\cdot\dT^{1/2}+\dT^{1/4}\!\cdot\dT^{1/4}\bigr)
   \ =\ C\bigl(\dT^{5/8}+\dT^{1/2}\bigr)\ \le\ C\,\dT^{1/2},
\end{equation}
the dominant term being $a\eta\cdot\eta$ with
$a\,\eta^2=\dT^{1/4}\cdot\dT^{1/4}=\dT^{1/2}$. As a cross--check against the
heuristic $\gd\sim\dT/\tau$: at $\tau=\sqrt\dT$ one has
$\dT/\tau=\dT^{1/2}$, matching \eqref{eq:s-gd}. Likewise the closeness of the
hole-filled inset \eqref{eq:m-fillclose} reads
\begin{equation}\label{eq:s-close}
   \abs{\Om\triangle\Etil}\ \le\ 8\pi\,\dT^{1/2}+A_0\,\dT^{1/2}
       +\frac{A_0^2}{16\pi^3}\,\dT+CH\ \le\ C\,\dT^{1/2},
\end{equation}
i.e.\ both contributions $\tau=\dT^{1/2}$ (the level height) and
$\dT/\tau=\dT^{1/2}$ (the level smoothing) coincide at the optimum, while the
spillover $\dT^2/\tau^2=\dT$ and the filled mass $H\le C\dT$ are genuinely
higher order. This is precisely the
balance $\tau\sim\dT/\tau$, solved by $\tau=\sqrt\dT$; any other power of
$\tau$ would leave one of the two terms larger than $\dT^{1/2}$.

\emph{Step 3: realisation as the measurable core graph $G$.} We take $G:=G_0$
the measurable open star-shaped core graph of \cref{prop:realization}, with
profile $\rho:=\rho_0\colon S^1\to[\rho_i,\rho_e]$ measurable. No
mollification or continuity is required: the nearly-spherical estimate
\cref{thm:graphG} admits measurable profiles. By
\cref{prop:realization}\ref{it:real-close} and \cref{cor:gd},
\begin{equation}\label{eq:s-EtilG}
   \abs{\Etil\triangle G}\ =\ \abs{\Etil\setminus G_0}
   \ \le\ C\,\eta(\Exc+a\eta)\ \le\ C\,\dT^{1/2},
\end{equation}
the last step being \eqref{eq:s-gd}, and
$\norm{\rho/\hat r_G-1}_{L^\infty}\le2\eta$ by
\cref{prop:realization}\ref{it:real-amp} with $\hat r_G=\hat r_{G_0}=
\sqrt{\abs G/\pi}\in[\rho_i,\rho_e]$.

\emph{Step 4: verification of (i), (ii), (iii).}

\noindent\textup{(i)} \emph{Amplitude.} By
\cref{prop:realization}\ref{it:real-amp} the profile $\rho=\rho_0$ obeys
$\norm{\rho/\hat r_G-1}_\infty\le2\eta$ with $\hat r_G=\sqrt{\abs G/\pi}\in
[\rho_i,\rho_e]$, and by \eqref{eq:s-eta} $\eta\le C\dT^{1/8}$, so
$\norm{\rho/\hat r_G-1}_\infty\le C\dT^{1/8}$. Shrinking $\delta_\star$ so
that $C\,\delta_\star^{1/8}\le\eta_G$ gives
$\norm{\rho/\hat r_G-1}_\infty\le\eta_G$, which is (i). Here it is essential
that $\eta\sim\dT^{1/8}$ is an \emph{absolute} smallness (decaying with
$\dT$), so the hypothesis of \cref{thm:graphG} is met for $\dT$ small,
\emph{regardless} of the cruder $L^1$--roundness $a\sim\dT^{1/4}$.

\noindent\textup{(iii)} \emph{Closeness.} By \eqref{eq:s-close} and
\eqref{eq:s-EtilG}, with $\Etil=E^h$ so that $\abs{\Om\triangle\Etil}$ already
includes the filled mass $H\le C\dT$ (\eqref{eq:m-fillclose}),
\[
   \abs{\Om\triangle G}\ \le\ \abs{\Om\triangle\Etil_0}+H+\abs{\Etil\triangle G}
   \ =\ \abs{\Om\triangle\Etil}+\abs{\Etil\triangle G}
   \ \le\ C\,\dT^{1/2}+C\,\dT^{1/2}\ =\ C\,\sqrt\dT,
\]
which is (iii); the filled mass $H\le C\dT$ is negligible.

\noindent\textup{(ii)} \emph{Deficit.} Here $G=G_0=
\{z+r\omega:r<\rho_0(\omega)\}$ is the measurable core graph of
\cref{cor:gd,prop:realization}, so $\rho_0\in
[\rho_i,\rho_e]$, $B_{\rho_i}(z)\subset G_0\subset\Etil\subset B_{\rho_e}(z)$,
and $\Etil\setminus G_0$ (the radial caps) lies in the collar
$\{\rho_i<\abs{x-z}<\rho_e\}$. The torsion-loss decomposition
\eqref{eq:torsion-loss} of \cref{lem:torsion-stab} holds unconditionally,
\[
   T(\Etil)-T(G_0)=\int_{\Etil\setminus G_0}u_{\Etil}+\int_{G_0}h
   \ \le\ \tilde\eta\,g+\int_{G_0}h,
   \qquad g=\abs{\Etil\setminus G_0},
\]
with $\tilde\eta=\tfrac14(\rho_e^2-\rho_i^2)\le\eta\,\rho_e\le C\dT^{1/8}$ and
the cap term $\int_{\Etil\setminus G_0}u_{\Etil}\le\tilde\eta\,g$ rigorous. The
torsion density is an absolute constant bounded below,
\[
   \frac{T(\Etil)}{\abs\Etil}
   =\bigl(1-\dstar(\Etil)\bigr)\frac{\abs\Etil}{8\pi}
   \ \ge\ \bigl(1-\tfrac12\bigr)\frac{\pi/2}{8\pi}=\frac1{32},
\]
using $\dstar(\Etil)\le\tfrac12$ (from Step~5 below, equivalently
\cref{prop:insetcontrols}\,\eqref{eq:dstar-inset}) and $\abs\Etil\ge\pi/2$.
Under the boundary-layer control \eqref{eq:flux-hyp} of
\cref{gap:torsion-stab} --- the single residual analytic point, namely the
localisation $\int_{G_0}h\le\tilde\eta\,g$ of the harmonic boundary-layer
term to the cap region --- the loss-slope hypothesis \eqref{eq:cutcond} holds
with $\lambda=2\tilde\eta\le\tfrac1{32}\le T(\Etil)/\abs\Etil$
(\cref{rem:supply,rem:roomtospare}), so \cref{lem:cutdeficit} applies to
$G=G_0\subset\Etil$ and gives directly, with no further regularisation step
(since $G=G_0$),
\begin{equation}\label{eq:s-defchain}
   \dstar(G)\ =\ \dstar(G_0)\ \le\ \dstar(\Etil).
\end{equation}
This is the only place in the reduction where a residual analytic point is
invoked; everything else --- the exact decomposition, the cap term, and the
deficit algebra --- is unconditional.

\emph{Step 5: transfer of the deficit from $\Om$ to $\Etil$.} The deficit of
the raw component $\Etil_0$ is controlled by the inset control
\eqref{eq:dstar-inset} of \cref{prop:insetcontrols},
$\dstar(\Etil_0)\le(\pi^2/\mu_E^2)\,\dstar(\Om)\le C\,\dstar(\Om)$; and the
hole-filling raises it only by the higher-order term $CH/\mu_E\le C\dT$
(\cref{lem:fill}\ref{it:fill-scales}, eq.~\eqref{eq:fill-dstar}), so
\begin{equation}\label{eq:s-transfer}
   \dstar(\Etil)\ \le\ \dstar(\Etil_0)+\frac{CH}{\mu_E}
   \ \le\ \frac{\pi^2}{\mu_E^2}\,\dstar(\Om)+C\dT\ \le\ C\,\dstar(\Om).
\end{equation}
We stress that the transfer to $\Etil_0$ is established in \cref{sec:inset}
\emph{without} appeal to the cutting principle \cref{lem:cutdeficit}: the
passage from the super--level set $\{u>\hat v\}$ to its largest component
$\Etil_0$ is handled there by the exact additivity of torsion over the
(disjoint) level components and the Saint--Venant bound on the spillover
torsion (\cref{prop:insetcontrols}, eqs.~\eqref{eq:Tadd}--\eqref{eq:Tspill}),
both internal to \cref{sec:identity,sec:inset}; and the hole-filling step
\eqref{eq:fill-dstar} uses only domain monotonicity of torsion and the area
bound $H\le C\dT$, again internal to \cref{sec:inset}. In particular the
spillover components need \emph{not} lie in the trapping collar of $\Etil$ ---
so \cref{lem:cutdeficit}, whose hypothesis requires the excised set to sit in
a thin concentric collar, is \emph{not} invoked here, and there is no circular
dependence on \cref{sec:cut} (see \cref{rem:fwd}). With $\mu_E\ge\pi/2$
(Step~0) the prefactor is $\pi^2/\mu_E^2\le4$, so for $\delta_\star$ small
\eqref{eq:s-transfer} gives $\dstar(\Etil)\le 5\,\dstar(\Om)$; in particular
$\dstar(\Etil)\le\tfrac12$ once $\dstar(\Om)\le\tfrac1{10}$, which justifies
the value used in Step~4. Inserting \eqref{eq:s-transfer} into
\eqref{eq:s-defchain} and using $\dT=\tfrac\pi8\dstar(\Om)\le C\dstar(\Om)$,
\[
   \dstar(G)\ \le\ \dstar(\Etil)+C\dT\ \le\ C\dstar(\Om)+C\dstar(\Om)
   \ =\ C\,\dstar(\Om),
\]
which is (ii). This completes the proof.
\end{proof}

\subsection{The main theorem}

\begin{theorem}[Sharp quantitative Saint--Venant, selection--free]
\label{thm:main}
There is an absolute constant $C$ such that every bounded open set
$\Om\subset\R^2$ with $\abs\Om=\pi$ satisfies
\begin{equation}\label{eq:main}
   \AF(\Om)^2\ \le\ C\,\dT(\Om).
\end{equation}
\end{theorem}

\begin{proof}
\emph{Trivial regime.} If $\dT\ge\delta_\star$, then since $\abs\Om=\abs B=\pi$
for the optimal disc $B$ we have $\abs{\Om\triangle B}\le\abs\Om+\abs B=2\pi$,
so $\AF(\Om)\le2$ and
$\AF(\Om)^2\le4\le(4/\delta_\star)\,\dT$. Thus \eqref{eq:main} holds with
$C=4/\delta_\star$ in this regime.

\emph{Near--extremal regime.} Assume $\dT\le\delta_\star$ with $\delta_\star$
the threshold of \cref{thm:reduction}, and let $G=\{z+r\omega:r<\rho(\omega)\}$
be the radial graph it provides, with $\hat r_G=\sqrt{\abs G/\pi}$ and
amplitude $\norm{\rho/\hat r_G-1}_\infty\le C\dT^{1/8}\le\eta_G$.

\Cref{thm:graphG} now applies directly to the graph $G$ as it is produced by
\cref{thm:reduction}, radial about the annulus centre $z$, with no
barycentre recentring: its hypotheses are the volume normalisation
$\hat r_G=\sqrt{|G|/\pi}$ (built into \cref{thm:reduction}), the $L^\infty$
amplitude bound $\|\rho/\hat r_G-1\|_\infty\le\eta_G$ (part~(i)), and the
quadratic mode-one bound \eqref{eq:mode1hyp}, i.e.\ $\|P_1\varphi\|_{L^2}\le
K\|\varphi\|_{L^2}^2$ for the profile $\varphi=\rho/\hat r_G-1$ in polar
coordinates about $z$. The first two hold by \cref{thm:reduction}. The third
is supplied by the choice of centre, as we now record.

\begin{remark}[The mode-one bound at the $L^1$-optimal centre]%
\label{gap:barycentre}
The centre $z$ is the \emph{exact} $L^1$-optimal (Fraenkel) centre of the
inset $\Etil$ (\cref{lem:zinterior}), an interior point with a definite
inscribed ball. Exact $L^1$-optimality is a first-order extremality condition
on the comparison ball: for the optimal centre the directional derivative of
$p\mapsto|\Etil\triangle B_{\hat r}(z+p)|$ vanishes, which is the statement
that the translation mode of $\partial\Etil$ carries no first-order
$L^1$-content. This is precisely the mechanism behind the quadratic mode-one
bound \eqref{eq:mode1hyp}: it forces the first Fourier mode of the radial
profile to be of second order in the amplitude, exactly as the barycentre
condition does through \cref{rem:mode1}, but \emph{at the centre $z$ already
in hand}, with no recentring of the set. What remains to make this fully
quantitative for the measurable core graph $G$ (as opposed to $\Etil$) is the
passage of the first-order extremality from the $L^1$-content of $\partial
\Etil$ to the $L^2$ mode-one coefficient of the profile of $G$, with the
collar error $\abs{\Etil\triangle G}\le C\sqrt\dT$ of \cref{thm:reduction}
absorbed. We isolate this as the explicit scoped hypothesis:
\begin{equation}\label{eq:mode1-scoped}
   \|P_1\varphi\|_{L^2(S^1)}\ \le\ K\,\|\varphi\|_{L^2(S^1)}^2
   \qquad(\varphi=\rho/\hat r_G-1\text{ about the }L^1\text{-optimal }z),
\end{equation}
for an absolute $K$. Under \eqref{eq:mode1-scoped} the argument below is
unconditional. This replaces the earlier, false-as-stated requirement of a
barycentre recentring --- which need not preserve single-valuedness of the
graph under an $L^\infty$ amplitude bound alone --- by a precise condition at
the existing centre, and removes the recentring obstruction entirely from
\cref{thm:graphG}. By \cref{lem:shift-asym} the asymmetry side needs no
mode-one hypothesis at all, and by \cref{rem:mode1-residual} the bound
\eqref{eq:mode1-scoped} is equivalent, for the present purpose, to the quartic
translation estimate $\dstar(G)\ge c\,\|P_1\varphi\|_{L^2}^4$
\eqref{eq:quartic-residual} --- the flatness of the translation direction of
the torsional deficit; this is the precise form of the residual.
\end{remark}

\emph{Asymmetry of $G$.} By \cref{thm:graphG} (with the mode-one bound
\eqref{eq:mode1-scoped}) and \cref{thm:reduction}(ii),
\begin{equation}\label{eq:m-AG}
   \AF(G)^2\ \le\ C_G\,\dstar(G)\ \le\ C_G\cdot C\,\dstar(\Om)
   \ =\ C\,\dstar(\Om)\ =\ C\,\tfrac8\pi\,\dT\ \le\ C\,\dT,
\end{equation}
so $\AF(G)\le C\sqrt\dT$.

\emph{Transfer to $\Om$ by the triangle inequality.} Let $B_G=B(y,\hat r_G)$
be the disc of area $\abs G$ realising the asymmetry of $G$, i.e.\ centred at
the optimal $y$ so that $\abs{G\triangle B_G}=\AF(G)\,\abs G$, and let
$B=B(y,1)$ be the disc of area $\pi$ \emph{concentric} with $B_G$. Since
$B_G$ and $B$ are concentric discs of areas $\abs G$ and $\pi$, their
symmetric difference is the annulus of area
\begin{equation}\label{eq:m-discs}
   \abs{B_G\triangle B}\ =\ \bigl|\,\abs G-\pi\,\bigr|
   \ =\ \bigl|\,\abs G-\abs\Om\,\bigr|\ \le\ \abs{\Om\triangle G}.
\end{equation}
Because $\abs\Om=\pi=\abs B$, the disc $B$ is admissible in the definition of
$\AF(\Om)$, so by the triangle inequality for the symmetric difference,
\begin{align*}
   \AF(\Om)\ \le\ \frac{\abs{\Om\triangle B}}{\pi}
   &\ \le\ \frac{\abs{\Om\triangle G}+\abs{G\triangle B_G}
            +\abs{B_G\triangle B}}{\pi}\\
   &\ \le\ \frac1\pi\Bigl(\abs{\Om\triangle G}+\AF(G)\,\abs G
            +\abs{\Om\triangle G}\Bigr),
\end{align*}
using \eqref{eq:m-discs} and $\abs{G\triangle B_G}=\AF(G)\abs G$. With
$\abs G=\pi+O(\sqrt\dT)\le2\pi$, $\abs{\Om\triangle G}\le C\sqrt\dT$ by
\cref{thm:reduction}(iii), and $\AF(G)\le C\sqrt\dT$ by \eqref{eq:m-AG},
\begin{equation}\label{eq:m-final}
   \AF(\Om)\ \le\ \frac1\pi\bigl(C\sqrt\dT+\AF(G)\,\abs G+C\sqrt\dT\bigr)
   \ \le\ C\sqrt\dT+C\,\AF(G)\ \le\ C\sqrt\dT.
\end{equation}
Squaring \eqref{eq:m-final} gives $\AF(\Om)^2\le C\,\dT$, which is
\eqref{eq:main} in the near--extremal regime. Combining the two regimes
proves the theorem.
\end{proof}

\begin{remark}[What was used, what was not, and the two sharp ingredients]
\label{rem:audit}
The reduction \cref{thm:reduction} draws on the two level deficits and
nothing else of substance:
\begin{itemize}[leftmargin=1.5em]
\item $D_H+D_I$ enters through the level identity
   $\int_0^m(D_H+D_I)\,dv=4\pi\dT$ (\cref{thm:level-identity}), which powers
   both the level selection \cref{lem:level} and the deficit transfer
   \cref{lem:transfer}, and through the pointwise bound $-\mu'\ge4\pi$
   (\cref{lem:muprime}) and the profile identity (\cref{lem:profile}) that
   carry closeness;
\item $D_I$ enters through the isoperimetric/annulus inputs
   (\cref{thm:annulus,thm:fmp}) and the multi--component spillover bound
   (\cref{lem:spill}).
\end{itemize}
At no point is the Brasco--De~Philippis--Velichkov selection principle used,
nor any mass--transport (moving--ball, kinetic) argument, nor a Serrin /
Magnanini--Poggesi symmetry estimate. The only nonlinear stability input is
the \emph{perturbative} nearly--spherical bound \cref{thm:graphG}, applied to
a graph whose amplitude is already absolutely small.

Two ingredients make the reduction \emph{sharp}, i.e.\ produce the exponent
$\tfrac12$ rather than $\tfrac14$:
\begin{enumerate}[label=\textup{(\alph*)},leftmargin=1.8em]
\item the fold--content bound \cref{prop:fold}, in which subtracting the
   constant $1/\hat r$ (legitimate because $\abs\Etil=\abs{B_{\hat r}(z)}$)
   and confining $\Etil\triangle B_{\hat r}(z)$ to the $\eta$--annulus
   replaces $a$ by $a\eta$, so the graph defect avoids the
   quantitative--isoperimetric square--root loss; without it the fold content
   would carry the full $a\sim\sqrt\diso$ and cap the result at
   $\dT^{1/4}$ closeness;
\item the cutting lemma \cref{lem:cutdeficit}: excising the low--torsion
   boundary caps to manufacture the graph $G$ \emph{decreases} the
   scale--invariant deficit $\dstar$, because the equal--area comparison ball
   shrinks faster (by $\gd/4$ in $\tilde\eta$--weighted terms) than the
   torsion is lost.
\end{enumerate}
Finally, the optimisation $\tau=\sqrt\dT$ is the balance between the two
$\tau$--dependent contributions to closeness: the level height $\sim\tau$
grows in $\tau$, while the smoothing/graph defect $\sim\dT/\tau$ decays in
$\tau$; equating them, $\tau\sim\dT/\tau$, forces $\tau=\sqrt\dT$, at which
both equal $\sqrt\dT$ and the amplitude is $\eta\sim(\dT/\tau)^{1/4}=
\dT^{1/8}$, the absolute smallness the graph theorem requires.
\end{remark}
